\subsection{Prior work on sequence distributions}\label{sec:relaxation}
Previously, RNA design was mostly approached with local search methods such as SAMFEO~\citep{zhou2023} and NEMO~\citep{portela2018}. These methods hold on to a few candidate sequences and improve them through small mutations. However, the design space grows exponentially with length, so these mutations can only reach a tiny part of it.

To address this, \citet{matthies2023} replaced the candidate sequences with a probability distribution over sequences, which is optimized with gradient descent. SamplingDesign~\citep{tang2026} then restricted this distribution, which they write as $\py(\cdot)$, to the valid design space $\GENX(\vecystar)$. Their goal is then to minimize the expected value of the objective function $f(\vecx,\vecystar)$ under the distribution.
\begin{equation}
\min_{\py}\;\obj,\quad\text{where }\obj\defeq\E_{\vecx\sim\py(\cdot)}\left[f(\vecx,\vecystar)\right].
\label{eq:minimize}
\end{equation}
They approximate this expected value with an average over a set of sampled sequences \SAMPLES.
\begin{equation}
\obj\approx\frac{1}{|\SAMPLES|}\sum_{\vecx\in\SAMPLES}f(\vecx,\vecystar).
\label{eq:expected}
\end{equation}
If the distribution puts all of its probability on one sequence \vecx, then \obj equals $f(\vecx,\vecystar)$, so this formula optimizes for that sequence only. The discrete problem in Eq.~\ref{eq:design} is therefore a special case of optimizing over a distribution.

\subsection{Coupled-variable distributions}\label{sec:variables}
\begin{figure}[t]
\centering
\begin{tikzpicture}[
  x=1.22cm,y=1.22cm,
  base/.style={circle,fill=white,line width=1.4pt,
    minimum size=0.42cm,inner sep=1pt,font=\normalsize\bfseries,text=black},
  position/.style={font=\normalsize\bfseries,text=black}
]
\filldraw[fill=vartwocolor!12,draw=vartwocolor!65,line width=0.8pt]
  (-1,0.865)--(1.431,1.365)--(1.431,0.365)--cycle;
\draw[varfourcolor!28,line width=9pt,line cap=round]
  (-0.809,2.681)--(0.809,2.681);
\draw[black!45,line width=1pt]
  (-0.5,0)--(-1,0.865)--(-0.5,1.73)--(-0.809,2.681)--
  (0,3.269)--(0.809,2.681)--(0.5,1.73)--
  (1.431,1.365)--(1.431,0.365)--(0.5,0);
\draw[varonecolor,line width=2.2pt] (-0.5,0)--(0.5,0);
\draw[varthreecolor,line width=2.2pt] (-0.5,1.73)--(0.5,1.73);
\node[font=\normalsize\bfseries,text=black,fill=white,inner sep=1pt]
  at (0,-0.32) {$v_1$};
\node[font=\normalsize\bfseries,text=black,fill=white,inner sep=1pt]
  at (0,1.46) {$v_3$};
\node[font=\normalsize\bfseries,text=black] at (0.66,0.85) {$v_2$};
\node[font=\normalsize\bfseries,text=black] at (0,2.681) {$v_4$};
\node[base,draw=varonecolor] (n1) at (-0.5,0) {G};
\node[base,draw=vartwocolor] (n2) at (-1,0.865) {A};
\node[base,draw=varthreecolor] (n3) at (-0.5,1.73) {C};
\node[base,draw=varfourcolor] (n4) at (-0.809,2.681) {C};
\node[base,draw=varfivecolor] (n5) at (0,3.269) {G};
\node[base,draw=varfourcolor] (n6) at (0.809,2.681) {A};
\node[base,draw=varthreecolor] (n7) at (0.5,1.73) {G};
\node[base,draw=vartwocolor] (n8) at (1.431,1.365) {U};
\node[base,draw=vartwocolor] (n9) at (1.431,0.365) {C};
\node[base,draw=varonecolor] (n10) at (0.5,0) {C};
\node[position] at (-0.5,-0.37) {1};
\node[position] at (-1.37,0.865) {2};
\node[position] at (-0.88,1.7) {3};
\node[position] at (-1.18,2.72) {4};
\node[position] at (0,3.64) {5};
\node[position] at (1.18,2.72) {6};
\node[position] at (0.84,1.98) {7};
\node[position] at (1.8,1.365) {8};
\node[position] at (1.8,0.365) {9};
\node[position] at (0.5,-0.37) {10};
\node[font=\normalsize\bfseries] at (-1.02,-0.08) {$5'$};
\node[font=\normalsize\bfseries] at (1.02,-0.08) {$3'$};
\end{tikzpicture}
\smallskip

{\normalsize
\begin{tabular}{@{}c@{\,}r@{\quad}l@{\quad}l@{}}
\textcolor{varonecolor}{\rule{0.65em}{0.65em}} & $v_1$ & Pair & $(1,10)$\\
\textcolor{vartwocolor}{\rule{0.65em}{0.65em}} & $v_2$ & Trimismatch & $(2,8,9)$\\
\textcolor{varthreecolor}{\rule{0.65em}{0.65em}} & $v_3$ & Pair & $(3,7)$\\
\textcolor{varfourcolor}{\rule{0.65em}{0.65em}} & $v_4$ & Mismatch & $(4,6)$\\
\textcolor{varfivecolor}{\rule{0.65em}{0.65em}} & $v_5$ & Unpaired & $(5)$
\end{tabular}}
\caption{Example of a target structure being split into variables. Each color is one variable, and the positions in a variable are chosen together.}
\label{fig:structureplot}
\end{figure}

To build this distribution, we turn the target structure into variables. Each variable is a group of positions whose nucleotides are chosen together, which we describe as coupling. We couple positions in the same way as SamplingDesign~\citep{tang2026}, which creates the four types of variables in Fig.~\ref{fig:structureplot}, namely unpaired positions, paired positions, mismatches, and trimismatches (1$\times$2 internal loops), with 4, 6, 16, and 64 choices. Paired positions are coupled so that every sampled pair is valid, and mismatches and trimismatches are coupled because the energy model scores these unpaired nucleotides together. Each variable has its own local distribution, and the probability of a sequence is the product of these local distributions.
\begin{equation}
\renewcommand{\arraystretch}{2.25}
\begin{array}{@{}r@{\;}>{\displaystyle}l@{}}
\py(\vecx)\defeq & \prod_{\mathrlap{\hspace{-0.75em}i\in\unpaired''(\vecystar)}}\;\pu_i(x_i)\\
\cdot & \prod_{\mathrlap{\hspace{-0.75em}(i,j)\in\pairs(\vecystar)}}\;\pp_{i,j}(x_ix_j)\\
\cdot & \prod_{\mathrlap{\hspace{-0.75em}(i,j)\in\mismatches(\vecystar)}}\;\pmm_{i,j}(x_ix_j)\\
\cdot & \prod_{\mathrlap{\hspace{-0.75em}(i,j,k)\in\trimismatches(\vecystar)}}\;\ptm_{i,j,k}(x_ix_jx_k).
\end{array}
\label{eq:factorization}
\end{equation}
Here $\unpaired''(\vecystar)$ holds the unpaired positions that are not in a mismatch or trimismatch. Any sequence drawn from this distribution is in $\GENX(\vecystar)$.

\subsection{SamplingDesign and its runtime bottleneck}\label{sec:samplingdesign}
SamplingDesign showed that optimizing this distribution produces high-quality designs, but this quality comes at a high computational cost, which is evident from the way it updates the distribution. SamplingDesign uses gradient descent, estimating the gradient of \obj from the sampled sequences at each step~\citep{tang2026}.
\begin{equation}
\nabla_{\py}\obj^{(t)}\approx\frac{1}{|\SAMPLES|}\sum_{\vecx\in\SAMPLES}\nabla_{\py}\log\py(\vecx)\,f(\vecx,\vecystar).
\label{eq:gradient}
\end{equation}
This estimated gradient is an average over all samples, so every sample must be scored at each step ($|\SAMPLES|=2500$ by default). Scoring a sample means running LinearPartition~\citep{zhang2020} on it, and this accounts for over 99\% of SamplingDesign's runtime~\citep{tang2026}, because scoring a large number of sequences is expensive even though LinearPartition itself is efficient. The gradient also only gives a direction, so the distribution moves a small amount at each step (a learning rate of 0.01 with up to 2,000 steps by default). As a result, SamplingDesign takes a long time to run, especially on longer puzzles, where a single step can take over 20 seconds~\citep{tang2026}.

In other words, the runtime of SamplingDesign depends heavily on the number of sequences scored, rather than on computing the gradient itself. This motivates the question of whether a distribution can learn from fewer samples and larger updates, possibly without a gradient at all.
